\subsubsection{RGB Images and Videos}\label{sec:rgb-images-and-videos}

Up to now we have represented a single image as:
\begin{equation*}
    I \in \mathbb{R}^{R \times C \times 3}
\end{equation*}
Where $R$ is the number of rows, $C$ is the number of columns, and $3$ is the number of channels (Red, Green, Blue). This representation is sufficient for photographs. But what about a \textbf{video}? A video is nothing more than a \hl{sequence of images.}

\highspace
Imagine recording one second of video. Instead of obtaining one image, we obtain many images taken one after another. Each image is called a \textbf{frame}. The number of frames per second is called the \textbf{frame rate} and is usually expressed in frames per second (fps). For example, a frame rate of 30 fps means that 30 images are recorded every second. Of course, the higher the frame rate, the smoother the video will be. For example, a video with a frame rate of 1 fps will look like a slideshow, while a video with a frame rate of 60 fps will look very smooth.

\highspace
Each individual frame is still:
\begin{equation*}
    I \in \mathbb{R}^{R \times C \times 3}
\end{equation*}
Nothing changes regarding the representation of a single image. However, since a video contains many frames, we simply add a \textbf{new dimension}. If there are $T$ frames, the video becomes:
\begin{equation*}
    V \in \mathbb{R}^{R \times C \times 3 \times T}
\end{equation*}
The first three dimensions $R \times C \times 3$ describe \textbf{space and color}, while the last dimension $T$ describes \textbf{time}. Roughly speaking, we ask each time: ``\emph{what color is this pixel \textbf{at this instant in time}?}''.

\highspace
\begin{flushleft}
    \textcolor{Green3}{\faIcon{question-circle} \textbf{Why does data size grow so quickly?}}
\end{flushleft}
The answer is very simple. In general, the total number of stored values is simply:
\begin{equation*}
    \text{Total number of values} = R \cdot C \cdot 3 \cdot T
\end{equation*}
For example, if we have a very small video of $R = 144$ rows, $C = 180$ columns, and $T = 30$ frames ($3$ channels), the total number of values is:
\begin{equation*}
    144 \cdot 180 \cdot 3 \cdot 30 = 2'332'800
\end{equation*}
About 2.3 million values! This is a very small video, but it is already quite large (on disk, it would take about $\approx 2$ MB of space).

\highspace
Every additional frame stores another complete image. If we double the number of frames, we double the required memory. If we double both image dimensions, memory grows by roughly four times. The increase is extremely fast. Let's imagine a full HD video with $R = 1080$ rows, $C = 1920$ columns, and $T = 30$ frames. The bytes required only for each frame are:
\begin{equation*}
    1080 \cdot 1920 \cdot 3 = 6'220'800 \approx 6 \text{ MB}
\end{equation*}
This is about 6.2 million values for each frame. If we have 30 frames, the total number of values is:
\begin{equation*}
    6'220'800 \cdot 30 = 186'624'000 \approx 187 \text{ MB}
\end{equation*}
This is about 187 million values per second!

\highspace
\textcolor{Green3}{\faIcon{question-circle} \textbf{With our phone, we can record videos at 60 fps, and we didn't see any explosions in memory. How is this possible?}} If one second requires roughly 187 MB, why is an MP4 often only a few megabytes? Because videos are \textbf{compressed}. Let's consider two consecutive frames of a video, and let's say that the first frame is a picture of a person standing in front of a wall. The second frame is the same person, but they have moved slightly to the right. The two frames are very similar, and the only difference is the position of the person. Instead of storing both frames completely, we can \hl{store only the first frame and then store only the difference between the two frames.} This is called \textbf{inter-frame compression}. Obviously there are many other compression techniques, but the main idea is to store only the information that changes between frames.

\highspace
\textcolor{Red2}{\faIcon{exclamation-triangle} \textbf{Important for Deep Learning.}} However, there is an important and unfortunate fact. Although videos are compressed on disk, \hl{neural networks do not process compressed files.} Before training, the \textbf{compressed video must be decoded into its full tensor}. This mans that a model actually receives a bomb of data like $\left(1080, 1920, 3, 30\right) = 186'624'000$ values. Thus, larger tensor imply more RAM, more GPU memory, slower training, large datasets, higher computational cost. This is one reason why video understanding is significantly more computationally demanding than image classification.